\documentclass[11pt, a4paper]{article}

% bfw-style.tex — gemeinsame Style-Definitionen für alle Handouts/Aufgabenhefte
% Voraussetzung: Kompilation mit xelatex (wegen fontspec)

\usepackage[ngerman]{babel}
\usepackage{geometry}
\geometry{a4paper, top=2.2cm, bottom=2.2cm, left=2.3cm, right=2.3cm}

\usepackage{fontspec}
% Fallback-Kette: Source Sans Pro -> Calibri -> Segoe UI -> Latin Modern Sans
% (Latin Modern ist immer mit MiKTeX vorhanden)
\IfFontExistsTF{Source Sans Pro}{%
  \setmainfont{Source Sans Pro}[Ligatures=TeX]%
}{\IfFontExistsTF{Calibri}{%
  \setmainfont{Calibri}[Ligatures=TeX]%
}{\IfFontExistsTF{Segoe UI}{%
  \setmainfont{Segoe UI}[Ligatures=TeX]%
}{%
  \setmainfont{Latin Modern Sans}[Ligatures=TeX]%
}}}
\IfFontExistsTF{Source Code Pro}{%
  \setmonofont{Source Code Pro}[Scale=0.92]%
}{\IfFontExistsTF{Consolas}{%
  \setmonofont{Consolas}[Scale=0.92]%
}{%
  \setmonofont{Latin Modern Mono}[Scale=0.92]%
}}

\usepackage{xcolor}
% Farbpalette: hell, freundlich, modern
\definecolor{bfwPrimary}{HTML}{0EA5A4}    % Petrol/Türkis - Primär
\definecolor{bfwPrimaryDark}{HTML}{0F766E} % dunkler Petrol für Headings
\definecolor{bfwAccent}{HTML}{F59E0B}     % Amber - Akzent
\definecolor{bfwSuccess}{HTML}{10B981}    % Grün - Erfolg / Lösung
\definecolor{bfwWarn}{HTML}{F97316}       % Orange - Achtung
\definecolor{bfwInk}{HTML}{0F172A}        % fast Schwarz
\definecolor{bfwInkSoft}{HTML}{475569}    % gedämpftes Grau
\definecolor{bfwBgSoft}{HTML}{F8FAFC}     % off-white
\definecolor{bfwBgTeal}{HTML}{ECFEFF}     % helles Türkis
\definecolor{bfwBgAmber}{HTML}{FEF3C7}    % helles Gelb
\definecolor{bfwBgRose}{HTML}{FFE4E6}     % helles Rosé für Eselsbrücken

\usepackage[colorlinks=true, linkcolor=bfwPrimaryDark, urlcolor=bfwPrimaryDark]{hyperref}
\usepackage{enumitem}
\setlist[itemize]{leftmargin=*, itemsep=2pt, topsep=2pt}
\setlist[enumerate]{leftmargin=*, itemsep=2pt, topsep=2pt}

\usepackage[most]{tcolorbox}
\usepackage{booktabs}
\usepackage{tikz}
\usetikzlibrary{decorations.pathmorphing, positioning, shapes.geometric, arrows.meta}

% Merkkästchen — wichtige Definition / Kernpunkt
\newtcolorbox{merksatz}[1][]{%
  enhanced, breakable,
  colback=bfwBgTeal,
  colframe=bfwPrimary,
  boxrule=0.6pt,
  arc=4pt,
  left=10pt, right=10pt, top=8pt, bottom=8pt,
  fonttitle=\bfseries\color{bfwPrimaryDark},
  title={\faLightbulb\ Merksatz},
  #1
}

% Eselsbrücke — Mnemoniks
\newtcolorbox{eselsbruecke}[1][]{%
  enhanced, breakable,
  colback=bfwBgRose,
  colframe=bfwAccent,
  boxrule=0.6pt,
  arc=4pt,
  left=10pt, right=10pt, top=8pt, bottom=8pt,
  fonttitle=\bfseries\color{bfwWarn},
  title={Eselsbrücke},
  #1
}

% Beispiel-Box
\newtcolorbox{beispiel}[1][]{%
  enhanced, breakable,
  colback=bfwBgSoft,
  colframe=bfwInkSoft,
  boxrule=0.4pt,
  arc=3pt,
  left=10pt, right=10pt, top=8pt, bottom=8pt,
  fonttitle=\bfseries\color{bfwInk},
  title={Beispiel},
  #1
}

% Lösungs-Box (für Aufgabenhefte)
\newtcolorbox{loesung}[1][]{%
  enhanced, breakable,
  colback=bfwBgAmber!40,
  colframe=bfwSuccess,
  boxrule=0.6pt,
  arc=3pt,
  left=10pt, right=10pt, top=8pt, bottom=8pt,
  fonttitle=\bfseries\color{bfwSuccess!70!black},
  title={Lösung},
  #1
}

% Glossar-Eintrag
\newcommand{\glossareintrag}[2]{%
  \par\noindent\textbf{\color{bfwPrimaryDark}#1} \enspace #2\par\smallskip
}

% Section-Stil — etwas farbiger als Standard
\usepackage{titlesec}
\titleformat{\section}
  {\Large\bfseries\color{bfwPrimaryDark}}
  {\thesection}{0.6em}{}
\titleformat{\subsection}
  {\large\bfseries\color{bfwPrimary}}
  {\thesubsection}{0.5em}{}
\titleformat{\subsubsection}
  {\normalsize\bfseries\color{bfwInk}}
  {\thesubsubsection}{0.4em}{}

% TikZ-Stil "skizzenhaft" — etwas weicher als Rough.js, aber stabil
\tikzset{
  roughbox/.style = {
    draw=bfwPrimaryDark, thick, fill=bfwBgTeal,
    rounded corners=4pt, inner sep=6pt
  },
  roughaccent/.style = {
    draw=bfwAccent, thick, fill=bfwBgAmber,
    rounded corners=4pt, inner sep=6pt
  },
  roughline/.style = {
    draw=bfwInkSoft, thick, -{Stealth[length=2.5mm]}
  }
}

% Bullet-Symbole (bewusst kein FontAwesome — vermeidet Font-Installations-Probleme)
\newcommand{\faLightbulb}{$\bullet$}
\newcommand{\faBookmark}{$\bullet$}

% Header/Footer
\usepackage{fancyhdr}
\pagestyle{fancy}
\fancyhf{}
\renewcommand{\headrulewidth}{0pt}
\fancyfoot[L]{\small\color{bfwInkSoft}\thepage}
\fancyfoot[R]{\small\color{bfwInkSoft} BFW M-064 Beschaffung im Büromanagement}


\title{\vspace*{-1.5cm}\Huge\bfseries\color{bfwPrimaryDark}Aufgabenheft Tag~10\\[0.4em]
  \large\color{bfwInkSoft}\textnormal{Postbearbeitung: Posteingang --- Übungen im IHK-Klausurformat}}
\author{\textsf{Büroprozesse gestalten und Arbeitsvorgänge organisieren \textperiodcentered{} M-067}\\
\textsf{\small\copyright\ Can Siebert\,\textperiodcentered\,2026}}
\date{Selbstlernmaterial \textperiodcentered{} 10.07.2026}

\begin{document}
\maketitle
\thispagestyle{empty}

\begin{merksatz}[title={\bfseries So nutzen Sie dieses Aufgabenheft}]
Bearbeiten Sie alle Aufgaben zunächst \emph{ohne} in die Lösungen zu schauen. Schreiben
Sie Ihre Antworten in vollständigen Sätzen --- die IHK-Prüfung verlangt formulierte
Begründungen, nicht bloße Stichworte. Beachten Sie die \emph{Operatoren} (nennen,
erläutern, beurteilen, begründen) --- sie geben den geforderten Antwort-Tiefegrad vor.
Erst danach öffnen Sie den Lösungsteil ab Seite~\pageref{loesungen} und vergleichen.
\end{merksatz}

\tableofcontents
\vspace{1em}
\hrule

\section{Aufgaben}

\subsection*{Aufgabe~1 --- Bedeutung des Posteingangs (10 Punkte)}

\noindent\textbf{\color{bfwAccent}YouTube-Vorbereitung:}\enspace \href{https://www.youtube.com/results?search_query=Posteingang+Poststelle+Buero}{Posteingang \& Poststelle} (\url{https://www.youtube.com/results?search_query=Posteingang+Poststelle+Buero})

\medskip
Die \emph{BüroKonzept KG} ist gewachsen und überlegt, eine eigene Poststelle einzurichten.

\begin{enumerate}[label=(\alph*)]
  \item \textbf{Erläutern} Sie, warum der Posteingang als „erster Umschlagplatz für Informationen” eine bedeutende Rolle spielt. \hfill (3~P.)
  \item \textbf{Nennen} Sie, wer in kleineren und wer in größeren Unternehmen die Tagespost bearbeitet. \hfill (2~P.)
  \item \textbf{Unterscheiden} Sie zentrale und dezentrale Poststelle und nennen Sie je einen Vorteil. \hfill (3~P.)
  \item \textbf{Beurteilen} Sie die Aussage: \emph{„Die Poststelle ist reine Hilfsarbeit ohne Bedeutung.”} \hfill (2~P.)
\end{enumerate}

\subsection*{Aufgabe~2 --- Möglichkeiten des Posteingangs (12 Punkte)}

\noindent\textbf{\color{bfwAccent}YouTube-Vorbereitung:}\enspace \href{https://www.youtube.com/results?search_query=Postzustellung+Briefmonopol+private+Dienstleister}{Postzustellung \& private Dienstleister} (\url{https://www.youtube.com/results?search_query=Postzustellung+Briefmonopol+private+Dienstleister})

\medskip
Die Post kann auf verschiedene Arten ein Unternehmen erreichen.

\begin{enumerate}[label=(\alph*)]
  \item \textbf{Nennen} Sie mindestens fünf Wege, auf denen Post in ein Unternehmen gelangen kann. \hfill (3~P.)
  \item \textbf{Erläutern} Sie, was ein \emph{Postfach} ist und wofür der \emph{Service Hin + Weg} dient. \hfill (3~P.)
  \item \textbf{Erklären} Sie, seit wann der deutsche Briefmarkt liberalisiert ist und welche Voraussetzung private Dienstleister benötigen. \hfill (3~P.)
  \item \textbf{Nennen} Sie zwei bekannte internationale private Brief- und Paketdienstleister. \hfill (3~P.)
\end{enumerate}

\subsection*{Aufgabe~3 --- Postvollmacht \& Arbeitsschritte (12 Punkte)}

\noindent\textbf{\color{bfwAccent}YouTube-Vorbereitung:}\enspace \href{https://www.youtube.com/results?search_query=Postvollmacht+Posteingang+Arbeitsschritte}{Postvollmacht \& Arbeitsschritte} (\url{https://www.youtube.com/results?search_query=Postvollmacht+Posteingang+Arbeitsschritte})

\medskip
Damit die Geschäftspost entgegengenommen werden darf, ist eine Postvollmacht nötig.

\begin{enumerate}[label=(\alph*)]
  \item \textbf{Erklären} Sie, was eine Postvollmacht ist und wozu sie benötigt wird. \hfill (3~P.)
  \item \textbf{Nennen} Sie vier Angaben, die eine Postvollmacht enthalten muss. \hfill (3~P.)
  \item \textbf{Bringen} Sie die Arbeitsschritte beim Posteingang in die richtige Reihenfolge: \emph{Stempeln, Annehmen, Verteilen, Öffnen, Vorsortieren, Kontrollieren}. \hfill (3~P.)
  \item \textbf{Erläutern} Sie, was mit einem \emph{Irrläufer} geschieht. \hfill (3~P.)
\end{enumerate}

\subsection*{Aufgabe~4 --- Post öffnen: darf oder darf nicht? (15 Punkte)}

\noindent\textbf{\color{bfwAccent}YouTube-Vorbereitung:}\enspace \href{https://www.youtube.com/results?search_query=Briefgeheimnis+persoenlich+vertraulich+Post+oeffnen}{Briefgeheimnis \& Post öffnen} (\url{https://www.youtube.com/results?search_query=Briefgeheimnis+persoenlich+vertraulich+Post+oeffnen})

\medskip
In der Poststelle der \emph{BüroKonzept KG} liegen vier Sendungen:

\begin{enumerate}[label=\arabic*., leftmargin=*]
  \item BüroKonzept KG / Herrn Achim Langner / Rochusstr.\ 30 / 53123 Bonn
  \item Persönlich / Frau Eva Holländer / BüroKonzept KG / Rochusstr.\ 30 / 53123 Bonn
  \item Herrn Sebastian Falo / BüroKonzept KG / Rochusstr.\ 30 / 53123 Bonn (ohne Zusatz)
  \item Büro GmbH / Herrn Volker Berg / Rochusstr.\ 320 (Empfänger dort unbekannt)
\end{enumerate}

\begin{enumerate}[label=(\alph*)]
  \item \textbf{Beurteilen} Sie für jede der vier Sendungen, ob sie geöffnet werden darf. \hfill (8~P.)
  \item \textbf{Erläutern} Sie, woran man im Anschriftenfeld Privatpost von Geschäftspost unterscheidet. \hfill (4~P.)
  \item \textbf{Beschreiben} Sie, was zu tun ist, wenn eine als „persönlich” gekennzeichnete Post versehentlich geöffnet wurde. \hfill (3~P.)
\end{enumerate}

\subsection*{Aufgabe~5 --- Kontrollieren \& Stempeln (12 Punkte)}

\noindent\textbf{\color{bfwAccent}YouTube-Vorbereitung:}\enspace \href{https://www.youtube.com/results?search_query=Posteingangsstempel+Leerkontrolle+Anlagenkontrolle}{Stempeln \& Kontrollieren} (\url{https://www.youtube.com/results?search_query=Posteingangsstempel+Leerkontrolle+Anlagenkontrolle})

\medskip
Nach dem Öffnen der Geschäftspost folgen Kontrolle und Stempelung.

\begin{enumerate}[label=(\alph*)]
  \item \textbf{Erläutern} Sie den Unterschied zwischen \emph{Leerkontrolle} und \emph{Anlagenkontrolle}. \hfill (3~P.)
  \item \textbf{Beschreiben} Sie, wo der Posteingangsstempel angebracht wird und welche Angaben er enthält. \hfill (3~P.)
  \item \textbf{Nennen} Sie drei Dokumente, die \emph{nicht} gestempelt werden, und erklären Sie, wie man stattdessen vorgeht. \hfill (3~P.)
  \item \textbf{Erläutern} Sie zwei Zwecke des Posteingangsstempels. \hfill (3~P.)
\end{enumerate}

\subsection*{Aufgabe~6 --- Verteilen, Digitalisierung \& Wertsendungen (12 Punkte)}

\noindent\textbf{\color{bfwAccent}YouTube-Vorbereitung:}\enspace \href{https://www.youtube.com/results?search_query=digitaler+Posteingang+Scannen+OCR}{Digitaler Posteingang \& Scannen} (\url{https://www.youtube.com/results?search_query=digitaler+Posteingang+Scannen+OCR})

\medskip
Ein Unternehmen mit hohem Postaufkommen führt ein digitales Posteingangssystem ein.

\begin{enumerate}[label=(\alph*)]
  \item \textbf{Nennen} Sie drei Möglichkeiten, wie die gestempelte Eingangspost verteilt werden kann. \hfill (3~P.)
  \item \textbf{Beschreiben} Sie die drei automatisierten Arbeitsschritte eines Posteingangssystems mit digitaler Archivierung. \hfill (3~P.)
  \item \textbf{Nennen} Sie vier Vorteile der Digitalisierung des Posteingangs. \hfill (4~P.)
  \item \textbf{Erläutern} Sie, wie Wertsendungen (z.\,B.\ Einschreiben) im Posteingang behandelt werden sollten. \hfill (2~P.)
\end{enumerate}

\vspace{1em}
\noindent\textbf{Punkte gesamt: 73}

\newpage

\section{Lösungen}\label{loesungen}

\subsection*{Lösung Aufgabe~1}
\begin{loesung}
\textbf{(a)} Der Posteingang ist der erste Umschlagplatz für ein- und ausgehende
Informationen. Über ihn laufen geschäftlich wichtige Schriftstücke wie Aufträge,
Rechnungen, Verträge und Behördenschreiben. Arbeitet er sorgfältig, gelangt die richtige
Post schnell zur richtigen Person; Fehler oder Verzögerungen wirken sich auf das ganze
Unternehmen aus (versäumte Fristen, verlorene Dokumente).

\textbf{(b)} In kleineren Unternehmen erledigt meist das \emph{Sekretariat} die Tagespost;
in größeren Unternehmen übernimmt eine eigene Abteilung, die \emph{Poststelle}, diese Aufgaben.

\textbf{(c)} \emph{Zentrale Poststelle}: eine Stelle bearbeitet die gesamte Post (Vorteil:
einheitliche Bearbeitung, Spezialisierung, effiziente Technik). \emph{Dezentrale Poststelle}:
jede Abteilung bearbeitet ihre Post selbst (Vorteil: kurze Wege, hohe Sachnähe).

\textbf{(d)} Die Aussage ist falsch. Die Poststelle entscheidet, ob wichtige Informationen
rechtzeitig und an der richtigen Stelle ankommen; sie sichert Fristen, das Briefgeheimnis und
die Nachvollziehbarkeit --- das ist für den gesamten Betrieb von hoher Bedeutung.
\end{loesung}

\subsection*{Lösung Aufgabe~2}
\begin{loesung}
\textbf{(a)} Mögliche Wege (fünf von): Zustellung durch Briefzusteller/Postbote (Brief,
Päckchen, Paket), Abholung aus einem Postfach, Service Hin + Weg, private Brief- und
Paketdienstleister, Fax, E-Mail, persönliche Abgabe am Empfang, Boten/Kuriere.

\textbf{(b)} Ein \emph{Postfach} wird in einer Postfiliale gemietet (einmalige
Einrichtungspauschale; Größe je nach Postaufkommen) und macht das Unternehmen unabhängig von
den Zustellzeiten des Postboten. Beim \emph{Service Hin + Weg} bringt die Deutsche Post AG die
Briefe zu vereinbarten Zeiten ins Unternehmen und holt die ausgehenden Briefe dort ab.

\textbf{(c)} Der deutsche Briefmarkt ist seit dem \emph{1.\ Januar 2008} liberalisiert
(Aufhebung des Briefmonopols). Private Dienstleister benötigen eine \emph{Lizenz der
Bundesnetzagentur}.

\textbf{(d)} Bekannte internationale Anbieter: z.\,B.\ TNT, FedEx, UPS.
\end{loesung}

\subsection*{Lösung Aufgabe~3}
\begin{loesung}
\textbf{(a)} Die Postvollmacht ist eine rechtsgültige, schriftliche Erklärung, mit der die
Geschäftspost (insbesondere Einschreiben und eigenhändige Sendungen) entgegengenommen werden
darf.

\textbf{(b)} Angaben (vier von): Vollmachtgeber, Bevollmächtigter, Gültigkeitsdauer, Umfang
der Vollmacht (normale Post, Einschreiben, eigenhändig), Datum der Ausstellung, Unterschrift
des Vollmachtgebers.

\textbf{(c)} Richtige Reihenfolge: \emph{Annehmen $\rightarrow$ Vorsortieren $\rightarrow$
Öffnen $\rightarrow$ Kontrollieren $\rightarrow$ Stempeln $\rightarrow$ Verteilen}.

\textbf{(d)} Ein Irrläufer ist eine fehlgeleitete Sendung. Sie wird \emph{ungeöffnet} dem
Postzusteller oder dem privaten Dienstleister zurückgegeben.
\end{loesung}

\subsection*{Lösung Aufgabe~4}
\begin{loesung}
\textbf{(a)} Beurteilung der vier Sendungen:
\begin{enumerate}[label=\arabic*., itemsep=0pt]
  \item Geschäftsbrief --- die Firma steht zuerst, dann der Name $\Rightarrow$ darf geöffnet und an Herrn Langner weitergeleitet werden.
  \item Privatbrief mit Zusatz „Persönlich” $\Rightarrow$ darf \emph{nicht} geöffnet werden, ungeöffnet weiterleiten.
  \item Name der Privatperson zuerst, Firma danach, \emph{kein} Zusatz „Persönlich/Vertraulich” $\Rightarrow$ darf laut Urteil (LAG Hamm 2003, AZ:~14 Sa 1972/02) in der Poststelle geöffnet werden.
  \item Irrläufer (falsche Adresse, Empfänger unbekannt) $\Rightarrow$ ungeöffnet dem Zusteller/Dienstleister zurückgeben.
\end{enumerate}

\textbf{(b)} \emph{Privatpost}: im Anschriftenfeld steht zuerst der Name einer natürlichen
Person, dann die Firma, oft mit Zusatz „Persönlich”/„Vertraulich”. \emph{Geschäftspost}:
zuerst die Firma, dann der Name.

\textbf{(c)} Den Brief wieder verschließen und mit dem Vermerk „versehentlich geöffnet”
weiterleiten.
\end{loesung}

\subsection*{Lösung Aufgabe~5}
\begin{loesung}
\textbf{(a)} \emph{Leerkontrolle}: Prüfung, ob der Umschlag nach Entnahme des Inhalts
vollständig geleert wurde. \emph{Anlagenkontrolle}: Prüfung, ob alle im Brief genannten
Anlagen beigefügt sind. Fehlt eine Anlage, wird dies handschriftlich auf dem Brief vermerkt.

\textbf{(b)} Der Posteingangsstempel wird \emph{rechts oben neben die Anschrift des
Empfängers} gesetzt. Er enthält das aktuelle Datum, evtl.\ die Uhrzeit und Platz für
Bearbeitungs- und Erledigungsvermerke.

\textbf{(c)} Nicht gestempelt werden z.\,B.\ Urkunden, Zeugnisse, Verträge, Schecks,
Zeichnungen. Stattdessen wird der \emph{Umschlag} gestempelt und dem Schriftstück beigefügt.

\textbf{(d)} Zwecke: Abgleich zwischen Datum des Schriftstücks und Eingangsdatum (Prüfung
verzögerter Zustellung) sowie Feststellung der innerbetrieblichen Bearbeitungszeit.
\end{loesung}

\subsection*{Lösung Aufgabe~6}
\begin{loesung}
\textbf{(a)} Verteilmöglichkeiten: Ablage in Mitarbeiter-/Abteilungsfächer (Post wird
abgeholt); ein Bote übernimmt die Verteilung; automatische Verteilung über
Schriftgutförderanlagen, z.\,B.\ ein Rohrpostsystem.

\textbf{(b)} Drei Arbeitsschritte: (1) maschinelle Öffnung der Briefumschläge mit elektrischen
Brieföffnern; (2) automatische Entnahme des Inhalts und elektronische Leerkontrolle; (3)
Einlesen der Schriftstücke mit Scannern und Speicherung im unternehmensinternen EDV-System
(Zugriff über den eigenen PC, häufig mit OCR-Texterkennung).

\textbf{(c)} Vorteile (vier von): schnellerer Informationsfluss; keine Zeitverzögerung durch
aufwändige Verteilung (kürzere Bearbeitungszeit); Schriftstücke jederzeit und schnell
auffindbar; sinkender Papierverbrauch; höhere Auskunftsfähigkeit; die Post kann nicht mehr
verloren gehen; höhere Effizienz der internen Prozessabläufe.

\textbf{(d)} Wertsendungen werden gegen Quittung angenommen (für Einschreiben/eigenhändig ist
die Postvollmacht nötig), in einem Eingangs-/Postbuch erfasst und sicher verwahrt bzw.\ gegen
Nachweis an den Empfänger übergeben.
\end{loesung}

\end{document}
